\documentclass[10pt,a4paper]{amsart}
\usepackage{amsmath,amssymb,amsthm,graphicx}
\usepackage{verbatim}
\newenvironment{reference}{\comment}{\endcomment}
%\newenvironment{reference}{}{}
\newenvironment{slogan}{\comment}{\endcomment}
\newenvironment{history}{\comment}{\endcomment}

% For commutative diagrams we use Xy-pic
\usepackage[all]{xy}

% We use 2cell for 2-commutative diagrams.
\xyoption{2cell}
% Theorem environments.
%
\theoremstyle{plain}
\newtheorem{theorem}[subsection]{Theorem}
\newtheorem{proposition}[subsection]{Proposition}
\newtheorem{lemma}[subsection]{Lemma}

\theoremstyle{definition}
\newtheorem{definition}[subsection]{Definition}
\newtheorem{example}[subsection]{Example}
\newtheorem{exercise}[subsection]{Exercise}
\newtheorem{situation}[subsection]{Situation}

\theoremstyle{remark}
\newtheorem{remark}[subsection]{Remark}
\newtheorem{remarks}[subsection]{Remarks}

\numberwithin{equation}{subsection}

% Macros
%
\def\lim{\mathop{\mathrm{lim}}\nolimits}
\def\colim{\mathop{\mathrm{colim}}\nolimits}
\def\Spec{\mathop{\mathrm{Spec}}}
\def\Hom{\mathop{\mathrm{Hom}}\nolimits}
\def\Ext{\mathop{\mathrm{Ext}}\nolimits}
\def\SheafHom{\mathop{\mathcal{H}\!\mathit{om}}\nolimits}
\def\SheafExt{\mathop{\mathcal{E}\!\mathit{xt}}\nolimits}
\def\Sch{\mathit{Sch}}
\def\Mor{\mathop{\mathrm{Mor}}\nolimits}
\def\Ob{\mathop{\mathrm{Ob}}\nolimits}
\def\Sh{\mathop{\mathit{Sh}}\nolimits}
\def\NL{\mathop{N\!L}\nolimits}
\def\CH{\mathop{\mathrm{CH}}\nolimits}
\def\proetale{{pro\text{-}\acute{e}tale}}
\def\etale{{\acute{e}tale}}
\def\QCoh{\mathit{QCoh}}
\def\Ker{\mathop{\mathrm{Ker}}}
\def\Im{\mathop{\mathrm{Im}}}
\def\Coker{\mathop{\mathrm{Coker}}}
\def\Coim{\mathop{\mathrm{Coim}}}

% Boxtimes
%
\DeclareMathSymbol{\boxtimes}{\mathbin}{AMSa}{"02}

%
% Macros for moduli stacks/spaces
%
\def\QCohstack{\mathcal{QC}\!\mathit{oh}}
\def\Cohstack{\mathcal{C}\!\mathit{oh}}
\def\Spacesstack{\mathcal{S}\!\mathit{paces}}
\def\Quotfunctor{\mathrm{Quot}}
\def\Hilbfunctor{\mathrm{Hilb}}
\def\Curvesstack{\mathcal{C}\!\mathit{urves}}
\def\Polarizedstack{\mathcal{P}\!\mathit{olarized}}
\def\Complexesstack{\mathcal{C}\!\mathit{omplexes}}
% \Pic is the operator that assigns to X its picard group, usage \Pic(X)
% \Picardstack_{X/B} denotes the Picard stack of X over B
% \Picardfunctor_{X/B} denotes the Picard functor of X over B
\def\Pic{\mathop{\mathrm{Pic}}\nolimits}
\def\Picardstack{\mathcal{P}\!\mathit{ic}}
\def\Picardfunctor{\mathrm{Pic}}
\def\Deformationcategory{\mathcal{D}\!\mathit{ef}}
\usepackage{cleveref}
\crefname{originalreference}{original reference}{original references}
\Crefname{originalreference}{Original reference}{Original references}
\makeatletter
\newlabel{sites-equation-stalk}{{04EH}{source}{}{sourcepointer-1}{}}
\newlabel{sites-equation-stalk@cref}{{[originalreference][0][]04EH}{[1][0][]source}{}{}{}}
\makeatother
\crefname{originalreference}{original reference}{original references}
\Crefname{originalreference}{Original reference}{Original references}
\makeatletter
\newlabel{algebra-lemma-characterize-local-ring}{{00E9}{source}{}{sourcepointer-2}{}}
\newlabel{algebra-lemma-characterize-local-ring@cref}{{[originalreference][0][]00E9}{[1][0][]source}{}{}{}}
\makeatother
\crefname{originallemma}{original lemma}{original lemmas}
\Crefname{originallemma}{Original lemma}{Original lemmas}
\makeatletter
\newlabel{lemma-locally-ringed}{{04ES}{source}{}{sourcepointer-3}{}}
\newlabel{lemma-locally-ringed@cref}{{[originallemma][0][]04ES}{[1][0][]source}{}{}{}}
\makeatother
\title{The locally ringed condition on a site versus locality of the stalks at its points}
\author{The Stacks Project Authors}
\date{stacks-project:sites-modules.tex:local-snapshot-20261012; GNU FDL 1.2 or later}
\begin{document}
\maketitle
\noindent\textit{Editorial note.} All delivered mathematical text is the author's own, copied verbatim from the pinned Stacks Project chapter file sites-modules.tex; no statement or proof step was rewritten, reordered or removed.\par\medskip
\noindent\textit{Editorial note.} The counted span cites three external items (a sites equation and related results), recorded in assembly.external\_references with canonical tags as navigation pointers.\par\medskip
\noindent\textit{Allowed standard prerequisites (editorial).}\begin{itemize}
\item Description of stalks at a point of a site as a filtered colimit over the point functor. Stacks Project, Sites, Equation sites-equation-stalk (tag 00Y5). Supplies the triple representation of stalk elements on which the argument runs.
\item Points of a site and the notion of enough points. Stacks Project, Sites, section on points of topoi. Provides the point functor used in (1) to (2) and the hypothesis under which the converse holds.
\end{itemize}\medskip
\ifdefined\hypertarget\hypertarget{sourcepointer-1}{}\fi
\noindent\textit{Source reference.} Original reference 04EH: sources/stacks-chow/tags.txt, source lines 5763--5763. Cross-chapter standard prerequisite cited by the author inside the counted span; it is used as a pointer and its own proof is not part of this item.\par\medskip
\ifdefined\hypertarget\hypertarget{sourcepointer-2}{}\fi
\noindent\textit{Source reference.} Original reference 00E9: sources/stacks-chow/tags.txt, source lines 526--526. Cross-chapter standard prerequisite cited by the author inside the counted span; it is used as a pointer and its own proof is not part of this item.\par\medskip
\ifdefined\hypertarget\hypertarget{sourcepointer-3}{}\fi
\noindent\textit{Source reference.} Original lemma 04ES: sources/stacks-chow/tags.txt, source lines 5773--5773. Cross-chapter standard prerequisite cited by the author inside the counted span; it is used as a pointer and its own proof is not part of this item.\par\medskip
\begin{lemma}
\label{lemma-locally-ringed-stalk}
Let $(\mathcal{C}, \mathcal{O})$ be a ringed site.
Consider the following conditions
\begin{enumerate}
\item For every object $U$ of $\mathcal{C}$ and $f \in \mathcal{O}(U)$
there exists a covering $\{U_j \to U\}$ such that for each $j$
either $f|_{U_j}$ is invertible or $(1 - f)|_{U_j}$ is invertible.
\item For every point $p$ of $\mathcal{C}$ the stalk $\mathcal{O}_p$
is either the zero ring or a local ring.
\end{enumerate}
We always have (1) $\Rightarrow$ (2). If $\mathcal{C}$ has enough points
then (1) and (2) are equivalent.
\end{lemma}

\begin{proof}
Assume (1). Let $p$ be a point of $\mathcal{C}$ given by a functor
$u : \mathcal{C} \to \textit{Sets}$. Let
$f_p \in \mathcal{O}_p$. Since $\mathcal{O}_p$ is computed by
Sites, Equation (\ref{sites-equation-stalk})
we may represent $f_p$ by a triple
$(U, x, f)$ where $x \in U(U)$ and $f \in \mathcal{O}(U)$.
By assumption there exists a covering $\{U_i \to U\}$
such that for each $i$ either $f$ or $1 - f$ is invertible
on $U_i$. Because $u$ defines a point of the site we see that
for some $i$ there exists an $x_i \in u(U_i)$ which maps to
$x \in u(U)$. By the discussion surrounding
Sites, Equation (\ref{sites-equation-stalk})
we see that $(U, x, f)$ and $(U_i, x_i, f|_{U_i})$ define the
same element of $\mathcal{O}_p$. Hence we conclude that
either $f_p$ or $1 - f_p$ is invertible. Thus
$\mathcal{O}_p$ is a ring such that for every element $a$
either $a$ or $1 - a$ is invertible. This means that $\mathcal{O}_p$
is either zero or a local ring, see
Algebra, Lemma \ref{algebra-lemma-characterize-local-ring}.

\medskip\noindent
Assume (2) and assume that $\mathcal{C}$ has enough points.
Consider the map of sheaves of sets
$$
\mathcal{O} \times \mathcal{O} \amalg \mathcal{O} \times \mathcal{O}
\longrightarrow
\mathcal{O} \times \mathcal{O}
$$
of
Lemma \ref{lemma-locally-ringed} part (3). For any local ring $R$
the corresponding map
$(R \times R) \amalg (R \times R) \to R \times R$
is surjective, see for example
Algebra, Lemma \ref{algebra-lemma-characterize-local-ring}.
Since each $\mathcal{O}_p$ is a local ring or zero the map is
surjective on stalks. Hence, by our assumption that $\mathcal{C}$
has enough points it is surjective and we win.
\end{proof}

\end{document}
